\documentclass[a4paper,11pt]{article}

\usepackage[ngerman]{babel}
\usepackage[top=3cm,bottom=3cm,left=3cm,right=3cm]{geometry} %NICHT IN KEMIE2


%\usepackage{microtype} %schöneres typesetting  mit [stretch=10] für nur 1% anstelle von 2%
%\usepackage{lmodern} %Latin Modern FONT
%\usepackage{pifont} %schönere Font



\usepackage{amsmath}
\usepackage{nicefrac}
\usepackage{amssymb}
%\usepackage{mathtools}

\usepackage{titlesec} %Um Section, etc bearbeiten zu können
\usepackage{setspace}
\usepackage{enumitem} %enumerate
\usepackage{multicol}
\usepackage[many]{tcolorbox}
\usepackage{hyperref}
\usepackage{arydshln}

%\renewcommand{\ttdefault}{lmtt} %Schriftart wird geändert zu Latin Modern Typewriter


\hypersetup{hidelinks,
			colorlinks=true,
			linkcolor=black,
			citecolor=miudiutex,
			urlcolor=miugray2
			}

\setlength{\parindent}{0pt}
\setcounter{section}{0}


\titlespacing*{\section}{0pt}{20mm}{6mm}
\titlespacing*{\subsection}{0pt}{6mm}{3mm}
\titlespacing*{\subsubsection}{0pt}{3mm}{0mm}



%\definecolor{miudiutex}{HTML}{2f3d39}
\definecolor{miudiutex}{HTML}{1d2926}
\definecolor{miugray}{HTML}{b4cccf}
\definecolor{miugray2}{HTML}{4d7365}
\definecolor{red}{HTML}{cf1f5c}
\definecolor{green}{HTML}{00A36C}
\definecolor{delphinidin}{HTML}{315794}
\definecolor{pelargonidin}{HTML}{b33807}
\definecolor{cyanidin}{HTML}{ba0746}
\definecolor{petunidin}{HTML}{b56fed}



\raggedcolumns %wichtig für Multicolumns


%================================
% SYMBOLS
%================================

\newcommand{\RA}{$\rightarrow$~}
\newcommand{\RL}{$\rightleftharpoons$~}
\newcommand{\HARP}{\ding{226}~}
%$\xrightarrow{2}$ %Pfeil mit 2 drüber


%================================
% SHORT COMMANDS (BOXES ; ETC)
%================================


\newcommand{\notiz}[1]{\begin{notizz}{}{}\begin{center}
			{\color{miugray2!50!miudiutex}#1} \end{center}\end{notizz}}
\newcommand{\zit}[2]{\begin{zitat*}{#1}{}\small #2 \end{zitat*}}

\newcommand{\frage}[2]{\begin{qst*}{#1}{}\small #2 \end{qst*}}

\newcommand{\unten}{\end{center} \tcblower \begin{center}}


\newcommand*\circled[1]{\tikz[baseline=(char.base)]{
            \node[shape=circle,draw,inner sep=0.5pt,miudiutex] (char) {\tiny #1};}}
            
\newcommand{\miusquare}{{\color{miugray2!50}\scriptsize $\blacksquare$}~}
\newcommand{\miusquarp}{{\color{petunidin!50}\scriptsize $\blacktriangleright$}~}
\newcommand{\niu}{\item[\miusquare]}
\newcommand{\nip}{\item[\miusquarp]}

%%%%% ABSAETZE

\newcommand{\pps}{\par \vspace*{1mm}}
\newcommand{\pp}{\par \vspace*{3mm}}
\newcommand{\ppbig}{\par \vspace*{8mm}}
\newcommand{\pphuge}{\par \vspace*{10mm}}

%================================
% FRAGE BOX
%================================

\tcbuselibrary{theorems,skins,hooks}
\newtcbtheorem{qst}{Frage:}
{%
	enhanced,
	breakable,
	colback = gray!15!white,
	frame hidden,
	boxrule = 0sp,
	borderline west = {2.5pt}{0pt}{red!70},
	sharp corners,
	detach title,
	before upper = \tcbtitle\par\smallskip,
	coltitle = black,
	fonttitle = \bfseries\sffamily,
	description font = \mdseries,
	separator sign none,
	segmentation style={solid, gray!50!black},
}
{th}




%================================
% ZITAT BOX
%================================

\tcbuselibrary{theorems,skins,hooks}
\newtcbtheorem{zitat}{Zitat}
{%
	enhanced,
	breakable,
	colback = gray!15!white,
	frame hidden,
	boxrule = 0sp,
	borderline west = {2.5pt}{0pt}{black!70},
	sharp corners,
	detach title,
%	before upper = \tcbtitle\par\smallskip,
	coltitle = gray!50!black,
	fonttitle = \bfseries\sffamily,
	description font = \mdseries,
%	separator sign none,
	segmentation style={solid, gray!50!black},
}
{th}



%================================
% SIMPLE SCHWARZER RAHMEN BOX
%================================


\newcommand{\boxfr}[1]{
	\begin{tcolorbox}[
	drop shadow southeast,
	enhanced,
	colframe=black!50,
	colback=white,
	boxrule = 1pt, 
%	toprule=0pt, bottomrule=0pt,rightrule=0pt,leftrule=0pt,
	boxsep=5pt,
	arc=1pt,
	]
	\begin{center}
	#1
	\end{center}
	\end{tcolorbox}
}

\definecolor{miudiutex}{HTML}{1d2926}
\definecolor{miugray}{HTML}{b4cccf}
\definecolor{miugray2}{HTML}{4d7365}
\definecolor{white2}{HTML}{f5f7f8}
\definecolor{red}{HTML}{cf1f5c}
\definecolor{green}{HTML}{00A36C}
\definecolor{delphinidin}{HTML}{315794}
\definecolor{uniblau}{HTML}{004e9f}
\definecolor{unigelb}{HTML}{fcba00}
\definecolor{unigrau}{HTML}{909085}
\definecolor{pelargonidin}{HTML}{b33807}
\definecolor{cyanidin}{HTML}{ba0746}
\definecolor{paeonidin}{HTML}{d15ca6}
\definecolor{petunidin}{HTML}{b56fed}
\definecolor{malvidin}{HTML}{8a2d8c}

\newcommand{\metermilli}{\mskip3mu \text{mm}}
\newcommand{\cm}{\mskip3mu \text{cm}}
\newcommand{\m}{ \text{m}}
\newcommand{\lite}{ \text{l}}
\newcommand{\kg}{ \text{kg}}
\newcommand{\g}{ \text{g}}
\newcommand{\km}{ \text{km}}
\newcommand{\h}{ \text{h}}
\newcommand{\s}{ \text{s}}
\newcommand{\tonne}{ \text{t}}
\usepackage[utf8]{inputenc}
\usepackage[T1]{fontenc}
\usepackage{lmodern}

\usepackage[
    activate={true,nocompatibility},
    final,
    tracking=true,
    kerning=true,
    spacing=true,
    factor=1100,
    stretch=30,
    shrink=20
]{microtype}
\usepackage{pdfpages}
\usepackage{awesomebox}
\usepackage{marginnote}
\tcbuselibrary{documentation}
\usepackage{sidenotes}
\usepackage{import}
\usepackage{xifthen}
\usepackage{transparent}
\usepackage[table]{xcolor}


\newcommand{\abbicon}{%
  \protect\tikz[baseline=0.2mm,scale=0.8]{%
    % Das blaue Quadrat mit abgerundeten Ecken
    \fill[uniblau, rounded corners=1.2pt] (0,0) rectangle (0.8em, 0.8em);
    % Der weiße Kreis in der Mitte
    \fill[white] (0.4em, 0.4em) circle (0.12em);
  }%
} 
\usepackage[labelfont={bf},font={color=miudiutex,small},figurename={\abbicon $~$Abbildung}]{caption}


\newcommand{\bckgrnd}{
\AddToHookNext{shipout/background}{
\begin{tikzpicture}[remember picture, overlay]
 \draw[draw=none,fill=uniblau!70, shift={(current page.north west)}] (-3,0) rectangle (23,-3.75);
 \draw[draw=none,fill=miugray, shift={(current page.north east)}] (-7,0) rectangle (3,-3.75);
\end{tikzpicture}
}
}



\newcommand{\incfig}[2]{%
\def\svgwidth{#2\columnwidth}
\import{C://LaTeX-Files/pngs/}{#1.pdf_tex}
}

\renewcommand{\textbf}[1]{\textsf{{\bfseries {#1}}}}



\titleformat{\section}[hang]
		{\LARGE \color{uniblau!75!black} \sffamily \bfseries}
		{\thesection}
		{6mm}
		{}
		[]
\titleformat{\subsection}[hang]
		{\large \sffamily \bfseries \color{black}}
		{\thesubsection}
		{4mm}
		{{\color{miugray}$\blacksquare ~$}}

\titleformat{\subsubsection}[block]
		{\normalsize \bfseries \sffamily \color{miudiutex}}%
		{\normalsize \color{black} \thesubsubsection}
		{2mm}
		{{\color{miugray}$\blacksquare ~$}}
		[ \color{black}]
		
\titleformat{\paragraph}
		{\normalsize \bfseries \sffamily \color{black}}
		{\footnotesize \theparagraph}
		{1mm}
		{}

\pagestyle{empty}
\titlespacing*{\section}{0pt}{20mm}{12mm}
\titlespacing{\section}{0pt}{20mm}{12mm}

\titlespacing*{\subsection}{0pt}{14mm}{4mm}
\titlespacing{\subsection}{0pt}{14mm}{4mm}

\titlespacing*{\subsubsection}{0pt}{6mm}{3mm}
\titlespacing{\subsubsection}{0pt}{6mm}{3mm}

\titlespacing*{\paragraph}{0pt}{6mm}{2mm}
\titlespacing{\paragraph}{0pt}{6mm}{2mm}


\let\oldmarginfigure\marginfigure
\let\endoldmarginfigure\endmarginfigure

\let\oldmarginfigure\marginfigure
\let\endoldmarginfigure\endmarginfigure

\renewenvironment{marginfigure}[1][0pt] % [1] steht für 1 Argument, [0pt] ist der Standardwert
  {\oldmarginfigure[#1]\captionsetup{labelfont={sf,bf,color=uniblau!75!black},font={color=miudiutex,small}}}
  {\endoldmarginfigure}
  
  
  
  
\renewcommand{\labelitemi}{\color{uniblau!70}$\bullet$}
\renewcommand{\labelitemii}{\color{uniblau!70}$\cdot$}
\newcommand{\plmtsquare}{{\color{uniblau!70}$\blacksquare~$}}
\newcommand{\splmt}[1]{{\color{uniblau!75!black} \sffamily #1}}
\newcommand{\fplmt}[1]{{\color{uniblau!75!black} \sffamily \bfseries #1}}

\newcommand{\antwort}[1]{
\begin{tcolorbox}[
	enhanced,
	enforce breakable,
	standard jigsaw,
	boxrule=0.7pt,
	colframe=black,
	sharp corners,
	boxsep=5pt,
	title=\bfseries \sffamily Antwort,
	opacityback=0,
	coltitle=miudiutex,
	attach boxed title to top text left={yshift=-\tcboxedtitleheight/2},
	boxed title style={colback=white,colframe=white},
	bottomrule at break=0pt,
	toprule at break=0pt,
	pad at break=16mm,
]
	\begin{center}
	#1
	\end{center}
	\end{tcolorbox}
}




\rowcolors{2}{miugray!50}{miugray!50}
\arrayrulecolor{white} % Weiße Gitterlinien
\setlength{\arrayrulewidth}{1.5pt} % Etwas dickere Linien
\renewcommand{\arraystretch}{1.4} % Mehr Zeilenhöhe für bessere Lesbarkeit}


%\setlist[itemize]{itemsep=0mm}
\setlist[enumerate]{itemsep=0mm}




\newif\ifshowme
\showmetrue




\begin{document}
\newgeometry{top=1.8cm,bottom=4cm,left=1.5cm,right=7.5cm,marginparsep=-2.00cm,marginparwidth=4.75cm}


\bckgrnd


\begin{center}
\LARGE
\color{white}
\textbf{Übung 11} \par \vspace*{2mm}
\large \color{miugray} \textbf{Prozessbezogene Lebensmitteltechnologie}
\end{center}
\par 
\vspace*{6mm}


\color{miudiutex}
\section*{Säuren \& Verteilungskoeffizient}


\subsection*{Aufgabe 1}
Der pKS-Wert von Buttersäure beträgt 4,82. Berechnen Sie den Anteil an undissoziierter Säure
bei den pH-Werten 1, 2, 3, 4, 5 und 6 und 7.

\ifshowme
\antwort{
\footnotesize
\begin{align*}
K_S &= \frac{[H_3O^+] \cdot [A^-]}{[HA]}  \\
pK_S &= -\log K_S \\
pH &= -\log [H_3O^+]
\end{align*}
Henderson-Hasselbalch:
\begin{align*}
pK_S &= pH -\log(\frac{[A^-]}{HA}) \\
 &= -\log [H_3O^+] -\log(\frac{[A^-]}{HA})
\end{align*}
Umstellen:
\begin{align*}
pH - pK_S &= \log(\frac{[A^-]}{HA}) & | 10^x \\
10^{pH - pK_S} &= \frac{[A^-]}{HA} \\
	&~\hat{=}~ \frac{1 - [HA]}{HA} \\
	&~\hat{=}~ \frac{1}{HA} - \frac{[HA]}{HA} \\
	&~\hat{=}~ \frac{1}{HA} - 1 \\
\\
10^{pH - pK_S}	&~\hat{=}~ \frac{1}{HA} - 1 \\
10^{pH - pK_S} +1 &= \frac{1}{HA} \\ \\
HA &= \frac{1}{10^{pH - pK_S} + 1 }
\end{align*}
\pp
Nachdem Einsetzen von pK$_S$ und pH-Wert erhält man folgende Ergebnisse: \pp
pH 1 \RA 99,985\% \par pH 2 \RA 99,85\% \par pH 3 \RA 98,5\% \par pH 4 \RA 86,9\% \par pH 5 \RA 39,8\% \par pH 6 \RA 6,2\% \par pH 7 \RA 0,656\%
}
\else
\fi


\subsection*{Aufgabe 2}
Bei der Entwicklung einer Light Mayonnaise mit einem Fettgehalt von 35\% steht die Frage im
Raum, ob Sorbinsäure oder Benzoesäure zugesetzt werden soll. Der pH-Wert der Mayonnaise
liegt bei pH 4,4.
\begin{enumerate}[label=\alph*)]
\item Berechnen Sie, wie hoch jeweils der Anteil an dissoziierter Säure wäre. Der pKS-Wert
von Benzoesäure liegt bei 4,18 und der pK$_S$-Wert von Sorbinsäure bei 4,76.

\ifshowme
\antwort{

Säurekonstante nach Dissoziationsgrad $\alpha$ umstellen:
\begin{align*}
K &= \frac{[H^+] \cdot [A]}{[HA]} \\
K &= \frac{[H^+] \cdot \alpha}{1 - \alpha} \\
K \cdot (1-\alpha) & = [H^+] \cdot \alpha \\
K - \alpha \cdot K &= [H^+] \cdot \alpha \\
K &= [H^+] \cdot \alpha + \alpha \cdot K \\
K &= \alpha \cdot ([H^+] + K) \\
\alpha &= \frac{K}{[H^+] + K}
\end{align*}
Nun können wir einfach einsetzen mit $K= 10^{-pK_S}$ und $[H^+] = 10^{-pH}$ \pp

\textbf{Benzoesäure:}
\begin{align*}
\alpha &= \frac{10^{-4,18}}{10^{-4,4} + 10^{-4,18}} \\
&= 0,624
\end{align*}
Daher liegt der Anteil an undissoziierter Säure bei 1-0,624 = 0,3759 \RA 37,6 \%
\pp 
\textbf{Sorbinsäure:}
\begin{align*}
\alpha &= \frac{10^{-4,76}}{10^{-4,4} + 10^{-4,76}} \\
&= 0,303
\end{align*}
Daher liegt der Anteil an undissoziierter Säure bei 1-0,303 = 0,697 \RA 69,7 \%
}
\else
\fi

\item Welche Relevanz hat das Ergebnis für die Auswahl des Konservierungsstoffs? Welche
weiteren Aspekte sind bei der Auswahl von Konservierungsstoffen zu beachten?
\ifshowme
\antwort{
\begin{itemize}
\item Für konservierende Wirkung ist der Anteil an undissozierter Säure ausschlaggebend.
Für den vorliegenden pH-Wert wäre also die Sorbinsäure stärker konservierend wirksam.
\item Betrachtung der ADI-Werte/Höchstmengen
\item Wirkungsspektrum und ''Ziel-MO'
\item Möglichkeit des kombinierten Einsatzes
\item Sensorische Eigenschaften des Produktes
\item Verteilungskoeffizient Fett/Wasser
\end{itemize}

}
\else
\fi

%\pagebreak
\item Nehmen Sie an, dass 200 mg/kg der Säure undissoziiert in der Mayonnaise vorliegen.
Wie hoch ist der Anteil in der wässrigen Phase? Verwenden Sie zur Abschätzung den
Verteilungskoeffizient K$_{OW}$ der beiden Säuren. Für Benzoesäure beträgt er in der
Mayonnaise 4,5, für Sorbinsäure beträgt er 8,7.

\ifshowme
\antwort{
Gibt man einen Stoff in ein System zweier nicht mischbarer Phasen, so verteilt er sich durch Diffusion,
bis sich ein Gleichgewicht eingestellt hat. Das Gleichgewicht wird durch den Nernst’schen
Verteilungskoeffizienten beschrieben:
\boxfr{
\begin{align*}
K_N = \frac{c_{p1}}{c_{p2}}
\end{align*}
}
\pp
\begin{align*}
K= 4,5 &= \frac{\frac{200~\text{mg}-x}{350}}{\frac{x}{650~\text{g}}} \\
4,5 &= \frac{200~\text{mg}-x}{350~\text{g}} \cdot \frac{650~\text{g}}{x}  \\
4,5 &= \frac{650}{350} \cdot (\frac{200~\text{mg}}{x}-1) \\ 
4,5 \cdot \frac{350}{650} &= \frac{200~\text{mg}}{x} -1 \\
(4,5 \cdot \frac{350}{650}+1) &= \frac{200~\text{mg}}{x} \\
\frac{1}{4,5 \cdot \frac{350}{650} + 1} \cdot 200 ~\text{mg} &= x  \\
x &= 58,43  ~\text{mg}
\end{align*}

Bei gleicher Verfahrensweise erhält man für \textbf{Sorbinsäure} (Verteilungskoeffizient = 8,7) einen Anteil von 35 mg in der wässrigen Phase
}
\else
\fi



\end{enumerate}
\subsection*{Aufgabe 3}
Propansäure hat einen pK$_S$-Wert von 4,88. Berechnen Sie den pH-Wert einer 0,1 molaren
Propansäurelösung.

\ifshowme
\antwort{
Formel für schwache Säuren:
\begin{align*}
pH &= \frac{1}{2} \cdot (pK_S - \log(c)) \\
pH &= \frac{1}{2} \cdot (4,88 - \log(0,1)) = 2,9
\end{align*}


}
\else
\fi





\end{document}
